\documentclass[10pt,a4paper]{amsart}
\usepackage[margin=19mm]{geometry}
\usepackage{amsmath,amssymb,mathtools}
\usepackage[scr=rsfs,cal=euler]{mathalfa}
\usepackage{xfrac}
\usepackage[unicode,hidelinks,bookmarks=false]{hyperref}
\newcommand{\ie}{i.e.\ }
\newcommand{\cf}{cf.\ }
\newcommand{\resp}{respectively\ }
\newcommand{\Subsection}{\subsection}
\providecommand{\nobreakdash}{\nobreak\hbox{-}}
\setlength{\emergencystretch}{2em}
\multlinegap0pt
\usepackage{amssymb}
\newtheorem{theorem}{Theorem}
\newtheorem{proposition}{Proposition}
\newtheorem{lemma}{Lemma}
\newcommand{\Aut}{\operatorname{Aut}}
\newcommand{\Fhat}{\widehat{F}}
\newcommand{\Hom}{\operatorname{Hom}}
\newcommand{\acl}{\operatorname{acl}}
\newcommand{\rk}{\operatorname{rk}}
\title{Subgroup measures and profinite rigidity in free groups}
\author{Shrinit Singh}
\date{Source: arXiv:2602.12815v2 (CC BY 4.0)}
\begin{document}
\maketitle

\noindent\emph{Source and licence.} Subgroup measures and profinite rigidity in free groups. Version arXiv:2602.12815v2 (CC BY 4.0), distributed by arXiv under the Creative Commons Attribution 4.0 International licence, http://creativecommons.org/licenses/by/4.0/, as recorded in the arXiv licence metadata for this exact version and shown on the abstract page https://arxiv.org/abs/2602.12815v2. Full licence notice: \texttt{licenses/arxiv-2602.12815-CC-BY-4.0.txt}.\par\smallskip

\noindent\textit{Editorial scope.} Counted unit: the article's theorem thm:universal (source0.tex lines 301--313). The article's complete original proof of it is delivered with it (source0.tex lines 954--1009, one \texttt{proof} environment of 56 lines, which the article places away from the statement). No mathematics is written, repaired or re-derived: the statement and the proof are the article's own lines, in the article's own notation, and no retained line is substituted or re-flowed.  Retained uncounted as the source-local context the proof needs: lines 441--453; lines 590--593; lines 632--638; lines 673--679; lines 711--714; lines 734--736.  Nothing was omitted from any retained span, so the package carries no omission record.  No retained line contains a citation, so the wrapper carries no bibliography and no bibliography entry.  No retained line contains a figure environment, an image-inclusion command or drawing code, so no image is delivered and the package declares no asset dependency.  Editorial formatting only, in addition: an AMS \texttt{amsart} preamble that restates the article's own declarations and macros used by the retained lines, namely the environments \texttt{theorem}, \texttt{proposition}, \texttt{lemma} on separate counters, together with the standard packages those lines need.  The retained environments print in this document as Theorem 1, Proposition 1, Lemma 1 in order of appearance, whereas the article prints the same items with the numbering of its own full document.  The complete native source of the article as distributed in the arXiv e-print for this exact version is bundled unchanged as \texttt{sources/194584/source0.tex}.
\begin{proposition}[Profinite orbit principle]\label{prop:lerf}
Let $G$ be a finitely generated LERF group, let $H_1,H_2\le G$ be finitely
generated subgroups, and let $\alpha\colon H_1\to H_2$ be an isomorphism.  The
following are equivalent.
\begin{enumerate}
\item[(i)] $\alpha$ is a measure equivalence relative to $G$;
\item[(ii)] the induced continuous isomorphism
\[
  \widehat\alpha\colon\overline{H_1}\longrightarrow\overline{H_2}
\]
extends to an automorphism of $\widehat G$.
\end{enumerate}
\end{proposition}

\begin{lemma}\label{lem:acl-trans}
Let $K\le P\le F$ with $F$ free of finite rank, $P$ a free factor of $F$ and $K$
finitely generated. Then $\acl_F(K)=\acl_P(K)$.
\end{lemma}

\begin{lemma}\label{lem:smallest-factor}
Let $F$ be a finitely generated free group and $K\le F$ a non-trivial finitely
generated subgroup. If $\widehat L$ is a free profinite factor of $\Fhat$ with
$\overline K\le\widehat L$, then $\overline{\acl_F(K)}\le\widehat L$.
Consequently $\overline{\acl_F(K)}$ is the smallest free profinite factor of
$\Fhat$ containing $\overline K$.
\end{lemma}

\begin{lemma}\label{lem:measure-stable}
Let $F$ and $P$ be finitely generated free groups and let $K\le F$ be finitely
generated. Then, for every finite group $G$ and every $h\in\Hom(K,G)$,
\[
  \Pr\nolimits^{F}_{K,G}(h)=\Pr\nolimits^{F*P}_{K,G}(h).
\]
\end{lemma}

\begin{lemma}\label{lem:aut}
For every $\psi\in\Aut(G)$ and every $H\le G$, the restriction
$\psi|_H\colon H\to\psi(H)$ is a measure equivalence.
\end{lemma}

\begin{lemma}\label{lem:groupoid}
Measure equivalences are closed under composition and inversion.
\end{lemma}

\begin{theorem}\label{thm:universal}
Let $F$ and $P$ be finitely generated free groups and let $K\le F$ be finitely
generated. Then
\begin{enumerate}
\item[(i)] $K$ is profinitely rigid in $F$ if and only if it is profinitely
  rigid in $F*P$;
\item[(ii)] $K$ is setwise profinitely rigid in $F$ if and only if it is setwise
  profinitely rigid in $F*P$.
\end{enumerate}
Consequently, marked and setwise profinite rigidity are stable under free
extension, and either property for $K\le F$ is equivalent to the corresponding
property in $\acl_F(K)$.
\end{theorem}

\begin{proof}[Proof of Theorem~\ref{thm:universal}]
We prove (i) and (ii) simultaneously. Throughout the proof, ``rigid'' means
respectively marked or setwise profinite rigidity. In the marked case we retain
the chosen measure equivalence $\alpha$ and require each witnessing ambient
automorphism to restrict to $\alpha$; in the setwise case this restriction
clause is omitted.

We may assume $K\neq1$, since the trivial subgroup is rigid in both senses in
every ambient group. Put $A=\acl_F(K)$ and write $F=A*C$. Since
$F*P=A*(C*P)$, Lemma~\ref{lem:measure-stable} shows that the measures induced on
$K$ are unchanged by adjoining a free factor. It is therefore enough to prove
the theorem when $K$ is algebraic in $F$, that is, when $\acl_F(K)=F$: the
algebraic case applied to $K\le A$ with free complements $C$ and $C*P$ then
identifies rigidity in $A$, $F$, and $F*P$. In the algebraic case
Lemma~\ref{lem:acl-trans} gives $\acl_{F*P}(K)=F$.

\medskip\noindent\emph{From $F$ to $F*P$.}
Let $L\le F*P$ be finitely generated and choose a measure equivalence
$\alpha\colon K\to L$. By Proposition~\ref{prop:lerf}, there is
$\Phi\in\Aut(\widehat{F*P})$ with $\Phi(\overline K)=\overline L$. Since
$\overline F$ is, by Lemma~\ref{lem:smallest-factor}, the smallest free
profinite factor containing $\overline K$, we have
\[
  \Phi(\overline F)=\overline{\acl_{F*P}(L)}.
\]
Thus $\rk\acl_{F*P}(L)=\rk F$. If
$F*P=\acl_{F*P}(L)*P'$, then $\rk P'=\rk P$; hence there is
$\beta\in\Aut(F*P)$ with $\beta(\acl_{F*P}(L))=F$. Replacing
$(L,\alpha)$ by $(\beta(L),\beta\circ\alpha)$ preserves measure
equivalence by Lemmas~\ref{lem:aut} and~\ref{lem:groupoid}; composing the
final witnessing automorphism with $\beta^{-1}$ recovers the original pair.
Thus we may assume $L\le F$. Lemma~\ref{lem:measure-stable} then identifies the comparison of
$K$ and $L$ in $F*P$ with the same comparison in $F$. Rigidity of $K$ in $F$
therefore gives $\varphi\in\Aut(F)$ with $\varphi(K)=L$ and, in the marked
case, $\varphi|_K=\alpha$. Extending $\varphi$ by the identity on $P$ proves
rigidity of $K$ in $F*P$.

\medskip\noindent\emph{From $F*P$ to $F$.}
Let $L\le F$ be finitely generated and choose a measure equivalence
$\alpha\colon K\to L$ computed in $F$. By Lemma~\ref{lem:measure-stable}, it is
also a measure equivalence in $F*P$. Rigidity there gives
$\varphi\in\Aut(F*P)$ with $\varphi(K)=L$ and, in the marked case,
$\varphi|_K=\alpha$. Automorphisms preserve algebraic closure, so
\[
  \varphi(F)=\varphi(\acl_{F*P}(K))
  =\acl_{F*P}(L)=\acl_F(L)\le F,
\]
where the equality $\acl_{F*P}(L)=\acl_F(L)$ is Lemma~\ref{lem:acl-trans}. Since $\varphi(F)$ is
a free factor of $F$ with the same finite rank as $F$, it follows that
$\varphi(F)=F$. Hence $\varphi|_F\in\Aut(F)$ witnesses rigidity of $K$ in $F$
(and restricts to $\alpha$ in the marked case).

The initial reduction through $A=\acl_F(K)$ also shows that either rigidity
notion for $K\le F$ is equivalent to the corresponding notion in
$\acl_F(K)$.
\end{proof}

\end{document}
